% ------------------------------------------------------------------------------
% 7 Results IV -- Machine learning and Bayesian inference
% ------------------------------------------------------------------------------
\section{Results IV: Machine-Learning Prediction and Bayesian
Inference}\label{sec:res-ml}

\subsection{Cross-validated prediction performance}

Table~\ref{tab:ml} and Figure~\ref{fig:mlcv} compare the five prognostic
models over 25 cross-validation folds. The random survival forest
achieves the best discrimination (cross-validated $C = 0.646$, IPCW AUC
at five years $0.677$) and the lowest integrated Brier score ($0.217$
over years 1--5), ahead of the clinical Cox model ($C = 0.634$),
gradient boosting ($0.633$), XGBoost ($0.625$), and the elastic-net
Cox model ($0.620$). The spread among models is modest---about 0.026
in the $C$-index---and the honest uncertainty (fold-to-fold SD
$\approx 0.025$) cautions against declaring a definitive winner; what
the comparison really shows is that, with ten baseline covariates and
$\sim$300 events, flexible learners extract only a little more signal
than a well-specified clinical model.

\begin{table}[htbp]
\centering\small
\caption{Machine-learning prediction of the composite DFS endpoint
(888 complete cases; $5\times5$-fold stratified cross-validation).
Treatment excluded from all predictors. AUC at five years is IPCW
weighted; IBS is the integrated Brier score over years 1--5.}
\label{tab:ml}
\begin{tabular*}{\textwidth}{@{\extracolsep{\fill}} l c c c c}
\toprule
Model & CV $C$-index & CV AUC (5y) & CV IBS & Apparent $C$ \\
\midrule
Random survival forest         & \textbf{0.646} & \textbf{0.677} & \textbf{0.217} & 0.713 \\
Clinical Cox (main effects)    & 0.634 & 0.652 & 0.222 & 0.651 \\
Gradient boosting (Cox loss)   & 0.633 & 0.662 & 0.225 & 0.717 \\
XGBoost (Cox objective) & 0.625 & -- & -- & 0.741 \\
Elastic-net Cox                & 0.620 & 0.646 & 0.226 & 0.628 \\
\bottomrule
\end{tabular*}
\end{table}

\begin{figure}[htbp]
  \centering
  \includegraphics[width=0.95\textwidth]{fig_ml_cv_performance.pdf}
  \caption{Cross-validated discrimination by model. Left: distribution
  of the $C$-index across 25 folds. Right: IPCW AUC at five years with
  cross-validation standard errors.}
  \label{fig:mlcv}
\end{figure}

The apparent-versus-cross-validated contrast in Table~\ref{tab:ml}
quantifies overfitting: XGBoost's apparent $C$ of $0.741$ drops to
$0.625$ under honest validation, a 0.12 optimism gap, versus a 0.07
gap for the random forest and essentially none for the clinical Cox
model. This gap is the single most important column in the table---it
is the difference between the performance one would report from a
naive in-sample fit and the performance a new patient would actually
experience.

Calibration of the winning model (Figure~\ref{fig:mlcal}, right panel)
is satisfactory across the risk spectrum: within deciles of predicted
five-year risk, the mean predicted risk tracks the Kaplan--Meier
observed risk within $6.5$ percentage points, with a slight
over-prediction in the lowest decile ($27.6\%$ predicted vs.\ $21.4\%$
observed). The five-year time-ROC achieves an AUC of $0.67$
(Figure~\ref{fig:mlcal}, left panel). Permutation importance
(Figure~\ref{fig:mlimp}) is dominated by the positive lymph-node
count (mean $C$-index drop $0.117$ per permutation, an order of
magnitude above every other covariate), followed at a distance by
advanced local extent, female sex, and a long surgery-to-registration
interval---a ranking that mirrors the adjusted Cox effects of
Figure~\ref{fig:forest} and reinforces the clinical coherence of the
learned model.

\begin{figure}[htbp]
  \centering
  \includegraphics[width=0.95\textwidth]{fig_ml_calibration_auc.pdf}
  \caption{Discrimination and calibration of the random survival
  forest. Left: ROC curve for five-year event classification (AUC =
  0.673). Right: predicted versus Kaplan--Meier observed five-year
  event risk by decile of predicted risk.}
  \label{fig:mlcal}
\end{figure}

\begin{figure}[htbp]
  \centering
  \includegraphics[width=0.90\textwidth]{fig_ml_importance.pdf}
  \caption{Permutation importance for the random survival forest
  (mean decrease in $C$-index over 20 shuffles per covariate, with
  standard deviations).}
  \label{fig:mlimp}
\end{figure}

\subsection{Bayesian re-analysis: priors and posteriors}\label{sec:res-bayes}

The Weibull model converged cleanly under all three prior scenarios
(two chains, target acceptance $0.9$; $\hat R \approx 1.00$ throughout).
Table~\ref{tab:bayes} and Figure~\ref{fig:bayesprior} display the
central result: the posterior hazard ratio is essentially insensitive
to the prior, ranging only from $0.591$ (weakly-informative) to $0.610$
(skeptical). Under the skeptical prior---which places 95\% of its mass
between hazard ratios of about $0.44$ and $2.28$ and only 5\% below
$0.5$---the posterior mean is $0.610$ with 95\% credible interval
$0.485$--$0.757$, $\Pr(\mathrm{HR}<1) = 1.000$ to three decimals, and
$\Pr(\mathrm{HR}<0.8) = 0.994$. The data are simply strong enough to
overwhelm any of the three priors, which is itself an informative
finding: conclusions of substance do not require prior consensus.

\begin{table}[htbp]
\centering\footnotesize
\caption{Bayesian Weibull analysis of the primary contrast under three
prior specifications (2 chains; 4{,}000 posterior draws each), with the
frequentist adjusted estimate for comparison.}
\label{tab:bayes}
\begin{tabular*}{\textwidth}{@{\extracolsep{\fill}} l c c c c c}
\toprule
Prior & Mean (log-HR) & Post.\ HR & 95\% CrI & $\Pr(\mathrm{HR}{<}1)$ & $\Pr(\mathrm{HR}{<}0.8)$ \\
\midrule
Skeptical $\mathcal{N}(0,\,0.421^2)$          & $0$        & 0.610 & 0.485--0.757 & 1.000 & 0.994 \\
Weakly-inf.\ $\mathcal{N}(0,\,1)$       & $0$        & 0.591 & 0.466--0.738 & 1.000 & 0.997 \\
Enthusiastic $\mathcal{N}(\log 0.65,\,0.421^2)$ & $-0.431$ & 0.592 & 0.474--0.734 & 1.000 & 0.997 \\
Frequentist (adj.\ Cox, MI) [ref.] & -- & 0.621 & 0.496--0.776 & -- & -- \\
\bottomrule
\end{tabular*}
\end{table}

\begin{figure}[htbp]
  \centering
  \includegraphics[width=0.95\textwidth]{fig_bayes_priors.pdf}
  \caption{Prior sensitivity of the Bayesian treatment-effect
  estimate. Left: the three prior densities on the hazard-ratio
  scale. Right: the corresponding posterior densities, with the
  frequentist point estimate (dotted) for reference.}
  \label{fig:bayesprior}
\end{figure}

\subsection{Bayesian sequential monitoring}

Figure~\ref{fig:bayesseq} replays evidence accumulation under the
skeptical prior at four event-driven looks. By the 25\% look (79
events), the posterior mean hazard ratio is already $0.601$ with
$\Pr(\mathrm{HR}<1) = 0.998$ and a credible interval excluding 1
(upper bound $0.88$); from the 50\% look (158 events) onward the
probability is essentially 1. Under the common Bayesian monitoring rule of
stopping when $\Pr(\mathrm{HR}<1)$ exceeds $0.975$, this trial would
have stopped at the \emph{first quarter} of the evidence---even
earlier than the frequentist group-sequential re-enactment of
Section~\ref{sec:res-design}, which required 50\% of the information
and a much larger $Z$. The contrast is partly definitional (the
Bayesian posterior probability does not pay a multiplicity price for
looks) and partly prior-driven (the skeptical prior pulls toward the
null early on, making its rapid crossing more impressive); both
mechanisms deserve emphasis when Bayesian designs are proposed as
replacements for alpha-spending.

\begin{figure}[htbp]
  \centering
  \includegraphics[width=0.95\textwidth]{fig_bayes_sequential.pdf}
  \caption{Bayesian sequential monitoring under the skeptical prior.
  Left: posterior mean hazard ratio with 95\% credible intervals at
  each look. Right: the monitoring statistic $\Pr(\mathrm{HR}<1)$
  against the 0.975 stopping threshold.}
  \label{fig:bayesseq}
\end{figure}

\subsection{Model criticism}

The posterior predictive check (Figure~\ref{fig:bayesppc}) replicates
the control arm's Kaplan--Meier curve from 100 posterior draws: the
observed curve lies comfortably within the 95\% posterior predictive
band throughout follow-up, with no systematic deviation in the tail
where a poorly specified baseline hazard would typically betray
itself. Approximate leave-one-out cross-validation strongly prefers
the Weibull baseline over the exponential special case (ELPD
difference $-18.0$, all stacking weight on the Weibull), and the
posterior shape parameter $\alpha = 0.765$ (95\% CrI $0.695$--$0.837$)
confirms a genuinely decreasing baseline hazard---patients who survive
the first two years have progressively lower event rates, exactly as
clinical experience suggests and as the Kaplan--Meier shape in
Figure~\ref{fig:km} already indicates.

\begin{figure}[htbp]
  \centering
  \includegraphics[width=0.90\textwidth]{fig_bayes_ppc.pdf}
  \caption{Posterior predictive check for the control arm: replicated
  Kaplan--Meier curves from 100 posterior draws versus the observed
  curve.}
  \label{fig:bayesppc}
\end{figure}
